\documentclass[tikz,border=1.5mm]{standalone}
\input{src/_preamble.tex}
\begin{document}
\begin{tikzpicture}[x=1mm,y=1mm,
  cellbg/.style={fill=pale, draw=white, line width=1pt, minimum width=22.5mm, minimum height=15mm},
  co/.style={font=\tiny, text=ink, align=center},
  hd/.style={font=\tiny\bfseries, text=ink, anchor=south, align=center},
  rl/.style={font=\tiny\bfseries, text=ink, anchor=east, align=right},
]
  \foreach \i/\h in {0/模型优先\\（不造本体，见第 31 章）,1/本体与模型全栈,2/工业部署优先,3/制造与\\垂直整合优先,4/研究出身＋\\母公司场景} {
    \node[hd] at ({22+\i*22.5},53) {\h};
  }
  \foreach \j/\r in {0/美国,1/中国,2/欧洲、日本、韩国} {\node[rl] at (9.8,{44-\j*15}) {\r};}
  \foreach \i in {0,...,4} {\foreach \j in {0,1,2} {\node[cellbg] at ({22+\i*22.5},{44-\j*15}) {};}}
  % 美国
  \node[co] at (22,44) {Physical Intelligence\\Skild AI\\Generalist AI};
  \node[co] at (44.5,44) {Figure\\1X};
  \node[co] at (67,44) {Agility\\Apptronik};
  \node[co] at (89.5,44) {特斯拉};
  \node[co] at (112,44) {波士顿动力\\（现代汽车集团）};
  % 中国
  \node[co] at (22,29) {（各家多自研模型；\\宇树、智元开源）};
  \node[co] at (44.5,29) {智元\\银河通用等};
  \node[co] at (67,29) {优必选};
  \node[co] at (89.5,29) {宇树\\小鹏、小米等车企};
  \node[co] at (112,29) {—};
  % 欧日韩
  \node[co] at (22,14) {—};
  \node[co] at (44.5,14) {Neura Robotics};
  \node[co] at (67,14) {ANYbotics（四足）};
  \node[co] at (89.5,14) {—};
  \node[co] at (112,14) {三星（Rainbow\\Robotics）};
\end{tikzpicture}
\end{document}
